\begin{table}[!ht]
\centering
\caption{Testes de robustez e diagnóstico. Cada linha traz o que foi
verificado, o resultado e a providência tomada. Valores extraídos dos
mesmos csv que alimentam as demais tabelas.}
\label{tab:robustez}
\small
\begin{tabular}{p{4.1cm}p{5.3cm}p{4.9cm}}
\toprule
Verificação & Resultado & Providência \\
\midrule
Heterocedasticidade (Breusch--Pagan) & $R^2$ da regressão auxiliar: 0{,}066 e 0{,}067 & erro-padrão agrupado por UPA e bootstrap em blocos \\
\addlinespace[2pt]
Forma funcional (RESET) & ganho de $R^2$ com potências do predito: 0{,}0020 a 0{,}0060 & rejeição sem relevância prática; forma mantida \\
\addlinespace[2pt]
Multicolinearidade (VIF) & VIF de \texttt{negro} $=$ 1{,}36; nenhum preditor acima de 10 & sem colinearidade relevante; o coeficiente de interesse não é afetado \\
\addlinespace[2pt]
Estrutura multinível (teste de razão de verossimilhança) & ICC da UPA no modelo nulo $=$ 0{,}353 & nível de bairro justificado (limiar de 5\%) \\
\addlinespace[2pt]
Método de estimação (ML vs.\ REML) & $\hat\tau^2$ da UPA: 0{,}30182 (ML) e 0{,}30182 (REML) & coincidem com $N$ de milhões; ML usado para os testes aninhados \\
\addlinespace[2pt]
Confundidor não observado (Konfound/ITCV) & 98{,}4\% das estimativas teriam de ser viés para anular o achado & achado robusto a confundimento plausível \\
\addlinespace[2pt]
Confundidor não observado (E-value) & E-value entre 1{,}46 e 2{,}36 & um confundidor precisaria dessa força com raça e desfecho \\
\addlinespace[2pt]
Sobreajuste (treino vs.\ teste) & maior diferença de $R^2$: 0{,}0084 & partição 80/20 e validação cruzada $k$-\textit{fold} \\
\addlinespace[2pt]
Ganho do ML sobre o linear & $R^2$ de teste: MQO 0{,}587; RF 0{,}589; XGBoost 0{,}644 & mesmo treino e teste; o ganho real está no XGBoost \\
\addlinespace[2pt]
Peso amostral (V1028) no acesso e na decomposição & OR do acesso 0{,}820 $\to$ 0{,}808; preço da Oaxaca (A) 17{,}9\% $\to$ 15{,}3\% & logit com UF fixo e erro agrupado por UPA; conclusão inalterada \\
\addlinespace[2pt]
Validação cruzada ($k=5$) & $R^2$ médio entre \textit{folds}: 0{,}6435 & hiperparâmetros escolhidos em partição de validação e confirmados por CV \\
\addlinespace[2pt]
Balanceamento e suporte comum & maior diferença padronizada entre grupos: 1{,}17 & desequilíbrio concentrado no contexto de bairro, absorvido pelo nível 2 \\
\addlinespace[2pt]
Especificação do nível de estado & coeficiente racial com UF aleatória: −0{,}0629 & igual ao da UF como efeito fixo; resultado não depende da escolha \\
\addlinespace[2pt]
Vazamento das médias de contexto & médias do bairro e do estado calculadas sem a própria observação & evita que a raça do indivíduo entre no regressor que o descreve \\
\addlinespace[2pt]
Dependência intragrupo (Moulton) & regressores que variam no nível da UPA & erro-padrão agrupado por UPA em todos os modelos \\
\bottomrule
\end{tabular}
\par\smallskip\footnotesize Fonte: Resultados originais da pesquisa.
\end{table}
